% ------------------------------------------------------------------------------
% Chapter 1
% ------------------------------------------------------------------------------
\chapter{Introduction} % enter the name of the chapter here
\label{cha:ch1_introduction} % enter the chapter label here (for cross referencing)

\section{Introduction} 
\label{sec:Introduction} 
This chapter introduces the study and also provides the background for examining the role of IT leadership in achieving ROI from AI coding assistants. It outlines the research problem being addressed and discusses the significance of the topic. Furthermore the chapter also presents the research questions and objections of this study. The final section of this chapter will provide an overview of the research methodology and theoretical framework used to guide the study.

\section{Background and context} 
\label{sec:Background and context} 
Artificial intelligence (AI) is one of the structural factors that has a significant impact on productivity and competitive advantages \autocite{Gillespie2025}. As of 2023, generative AI systems have developed from being an experiment to reaching the phase of implementation within organizational processes \autocite{Gillespie2025}. Recent reports from around the world show that organizations allocate more money to AI since they expect clear monetary outcomes, cost savings, and improved performance of organizational processes \autocite{Gillespie2025}. Research on information systems and management has transcended from the stage of adoption to the understanding of added value meanings with the emphasis on differentiation of organizational value created from investments in AI systems as a function of organizational governance and management performance, not technology excellence \autocite{Hossain2025}. 

In software engineering, the use of generative AI tools are among the most noticeable forms of AI-enabled augmentation \autocite{Peng2023}. Providing developers with functionalities such as code suggestions and entire code functions in real time based on context, documentation, and debugging \autocite{Peng2023}.Practical controlled experiments provide conclusive evidence that AI-assisted developers complete software development tasks faster and report increased productivity of about 55.8\% compared with control groups \autocite{Peng2023}. Collaboration between humans and AI systems is a paradigm shift from the traditional concept of automation, in which AI enhances developers’ capabilities rather than replacing them with machines \autocite{Peng2023}. Initial performance improvements demonstrated in controlled promising experiments may not necessarily guarantee long-term organizational performance benefits. Empirical studies report considerable variation in the productivity gains from AI-assisted software development tasks across software development teams and organizations. 

As it has been long noted, organization culture and management style have been generally believed to affect the implementation success of new technologies. Within the digital sphere, values and practices will drive the culture of experimentation, learning, and data use in decision making as well as a positive attitude to tech adoption. New findings reveal that digital culture explains the linkage between human AI skills and innovation outcomes, and that culture readiness mediates the performance impact of AI Investments \autocite{An2024}. Similarly, this applied also to employees' engagement in AI, ethical AI management or leadership impacts toward strategic alignment. Even though the significance of digital culture and leadership is being gradually recognized by most researchers, current research findings on AI Adoption is more fragmented \autocite{An2024}. A significant number of recent works investigated AI adoption under the macro-organizational view or focused on the impacts for personal performance only instead of integrating digital culture, management and ROI in the software industry \autocite{Peng2023}.

Generative AI (GenAI) has been increasingly adopted by african companies to address resource constraints and operational inefficiencies. Similarly, there is evidence that developing nations have tried extensively with AI and have positive views about AI’s capabilities to boost productivity at work \autocite{Gillespie2025}. Yet, literature to date highlights difficulties around digital skills pipelines, governance maturity and leadership capability among others in adopting GenAI technologies \autocite{Hughes2026}. In this context, organisational culture and leadership remain particularly critical. If these factors are not supported by a clear strategic vision, trust and a learning orientation, AI adoption is likely to be less about value creation and more about segmentation or symbolism. The digital transformation landscape in Africa is a valid environment for exploring the impact of socio-organizational factors on the achievement of returns from AI technologies.



\section{Rationale of the study} 
\label{sec:Rationale of the study} 

\section{Problem statement} 
\label{sec:Problem statement} 

\section{Significance of the study} 
\label{sec:Significance of the study} 

\section{Overview of the theoretical/conceptual framework} 
\label{sec:Overview of the theoretical/conceptual framework} 

\section{Research aim and objectives} 
\label{sec:Research aim and objectives} 

\section{Research questions} 
\label{sec:Research questions} 

\section{Scope and delimitation} 
\label{sec:Scope and delimitation} 

\section{Methodology overview} 
\label{sec:Methodology overview} 

\section{Ethical considerations} 
\label{sec:Ethical considerations} 

\section{Dissertation structure} 
\label{sec:Dissertation structure} 

\section{Chapter summary} 
\label{sec:Chapter summary} 